\usepackage[T1]{fontenc}
\usepackage[utf8]{inputenc}
\usepackage{lmodern}
\usepackage{amsmath,amssymb,amsthm,mathtools,mathrsfs}
\usepackage[margin=1in,headheight=14pt]{geometry}
\usepackage{microtype}
\usepackage{booktabs,longtable,array,enumitem}
\usepackage{xcolor}
\definecolor{accent}{HTML}{184B68}
\definecolor{linkblue}{HTML}{155D8B}
\usepackage{fancyhdr}
\pagestyle{fancy}
\fancyhf{}
\fancyhead[L]{\small\sffamily Tail kernels and boundary products}
\fancyhead[R]{\small\sffamily ProveIt}
\fancyfoot[C]{\thepage}
\renewcommand{\headrulewidth}{0.3pt}
\usepackage{hyperref}
\hypersetup{colorlinks=true,linkcolor=linkblue,citecolor=linkblue,urlcolor=linkblue,
  pdftitle={Bilinear Tail Identities, Polylogarithmic Resonances, and Stieltjes Boundary Products},
  pdfauthor={ProveIt research continuation},bookmarksnumbered=true}
\usepackage{bookmark}
\numberwithin{equation}{section}
\newtheorem{theorem}{Theorem}[section]
\newtheorem{proposition}[theorem]{Proposition}
\newtheorem{lemma}[theorem]{Lemma}
\newtheorem{corollary}[theorem]{Corollary}
\theoremstyle{definition}
\newtheorem{definition}[theorem]{Definition}
\newtheorem{example}[theorem]{Example}
\theoremstyle{remark}
\newtheorem{remark}[theorem]{Remark}
\newtheorem{question}[theorem]{Question}
\DeclareMathOperator{\Li}{Li}
\DeclareMathOperator{\FP}{FP}
\DeclareMathOperator{\CT}{CT}
\DeclareMathOperator{\Res}{Res}
\DeclareMathOperator{\rank}{rank}
\DeclareMathOperator{\diag}{diag}
\DeclareMathOperator{\ord}{ord}
\newcommand{\C}{\mathbb C}
\newcommand{\Q}{\mathbb Q}
\newcommand{\N}{\mathbb N}
\newcommand{\Z}{\mathbb Z}
\newcommand{\Bcal}{\mathcal B}
\newcommand{\fall}[2]{#1^{\underline{#2}}}
\newcommand{\sourcepin}{da2792fc9cc8ef19c6a314be6c65e1437fbdc66b}
\setlength{\parskip}{3pt}
\setlength{\emergencystretch}{2em}
\setcounter{tocdepth}{2}
\allowdisplaybreaks[2]
\setlist{topsep=4pt,itemsep=3pt}
